\section{De VAE a VAE jerárquico}

\subsection{Motivación para agregar niveles}

Las limitaciones de un VAE radican en que una única distribución gaussiana resulta difícil de ajustar a distribuciones de datos complejas. Una idea natural es: dejar de usar una sola variable latente $z$ e introducir \textbf{variables latentes múltiples} $x_1, x_2, \ldots, x_T$, formando así una estructura jerárquica.

\begin{knowledgebox}{Intuición sobre cómo mejora la expresividad al agregar niveles}
Una distribución gaussiana individual solo cuenta con dos grupos de parámetros: media y varianza, lo que limita su capacidad expresiva. Pero si tomamos el producto de múltiples distribuciones de transición gaussiana como la verosimilitud global, entonces aunque cada paso de transición siga siendo una gaussiana simple, \textbf{la distribución global puede ser mucho más compleja que una única gaussiana}. Esto es análogo a cómo un camino recto simple puede aproximar cualquier curva compleja mediante múltiples pequeños ajustes de dirección.
\end{knowledgebox}

\subsection{Cambio de notación}

A partir de aquí, cambiamos el sistema de notación para adaptarlo al marco de los modelos de difusión:

\begin{itemize}
    \item $x_0$: punto de datos original (anteriormente $x$)
    \item $x_1, x_2, \ldots, x_T$: variables latentes múltiples (anteriormente $z$)
    \item $x_T$: ruido puro final, que sigue $\mathcal{N}(0, I)$
\end{itemize}

\begin{warningbox}{Recordatorio sobre el cambio de notación}
Tenga cuidado con el cambio de notación de VAE a modelos de difusión: antes se usaba $z$ para las variables latentes y $x$ para los datos; ahora se utiliza uniformemente $x_t$ para representar las variables en distintos pasos temporales. Esto se debe que, en los modelos de difusión, todas las variables intermedias y los datos poseen \textbf{las mismas dimensiones}, pues residen en el mismo espacio. Distintos artículos pueden emplear sistemas de notación diferentes; al leerlos, es necesario prestar atención a esto.
\end{warningbox}

\subsection{Hipótesis de Markov}

Para simplificar el entrenamiento del VAE jerárquico, introducimos la \textbf{hipótesis de Markov}: la muestreo de cada $x_t$ depende únicamente del paso anterior $x_{t-1}$, y no de estados anteriores más tempranos.

$$q(x_t | x_{t-1}, x_{t-2}, \ldots, x_0) = q(x_t | x_{t-1})$$

Bajo la hipótesis de Markov, toda la distribución conjunta puede descomponerse en el producto de distribuciones de transición:

\textbf{Proceso generativo} (verosimilitud/proceso inverso):
$$p_\theta(x_0, x_1, \ldots, x_T) = p(x_T) \prod_{t=1}^{T} p_\theta(x_{t-1} | x_t)$$

\textbf{Proceso de codificación} (variational posterior/proceso hacia adelante):
$$q(x_1, \ldots, x_T | x_0) = \prod_{t=1}^{T} q(x_t | x_{t-1})$$

\begin{importantbox}{Importancia de la hipótesis de Markov}
Sin asumir las propiedades de Markov, el muestreo de $x_t$ podría depender de todos los estados históricos $x_0, x_1, \ldots, x_{t-1}$, lo que haría que el entrenamiento fuera extremadamente complejo. La hipótesis de Markov simplifica enormemente los cálculos, permitiendo que ELBO sea descompuesto y optimizado eficazmente. En cursos posteriores se discutirán casos no markovianos (como DDIM).
\end{importantbox}

\subsection{Desafíos del VAE jerárquico}

Aunque la estructura jerárquica incrementa la capacidad expresiva, aún enfrenta desafíos bajo el marco estándar de VAE:

\begin{enumerate}
    \item \textbf{Dificultad en el entrenamiento conjunto}: se requiere entrenar simultáneamente el codificador (aprendiendo la variational posterior $q_\phi$) y el decodificador (aprendiendo la verosimilitud $p_\theta$); cuanto más niveles haya, menos estable será el entrenamiento.
    \item \textbf{Costo computacional elevado}: el término de consistencia en ELBO involucra múltiples variables aleatorias, lo que resulta costoso de calcular.
\end{enumerate}

Estos desafíos impulsaron las elecciones de diseño centrales de los modelos de difusión: \textbf{fijar el proceso hacia adelante y aprender únicamente el proceso inverso}.

\subsection{Resumen de la sección}

Mediante la introducción de variables latentes múltiples y la hipótesis de Markov, el VAE jerárquico mejora su capacidad expresiva. No obstante, entrenar conjuntamente codificador y decodificador sigue siendo difícil, lo que nos impulsa a buscar una solución que no requiera aprender un codificador: es decir, los modelos de difusión.



